\documentclass[11pt,a4paper]{article}
\usepackage[utf8]{inputenc}
\usepackage[T1]{fontenc}
\usepackage[brazil]{babel}
\usepackage{lmodern}
\usepackage{amsmath,amssymb,amsthm}
\usepackage{booktabs}
\usepackage{float}
\usepackage[margin=2.5cm]{geometry}
\usepackage{pgfplots}
\pgfplotsset{compat=1.16}
\usepackage[hidelinks]{hyperref}

\newtheorem{teorema}{Teorema}
\newtheorem{proposicao}[teorema]{Proposição}
\newtheorem{lema}[teorema]{Lema}
\theoremstyle{remark}
\newtheorem{observacao}[teorema]{Observação}

\title{Uma aproximação para o perímetro da elipse\\ via números de Catalan e polilogaritmos}
\author{}
\date{Novembro de 2026}

\begin{document}
\maketitle

\begin{abstract}
Os coeficientes da série de Gauss--Kummer para o perímetro da elipse são
quadrados de números de Catalan divididos por potências de 2. A assintótica clássica
dos números de Catalan mostra que esses coeficientes se comportam como
$\frac{1}{4\pi n^3}\big(1+\frac{3}{4n}+\cdots\big)$, ou seja, como os
coeficientes de uma combinação de polilogaritmos
$\mathrm{Li}_k(h)=\sum_n h^n/n^k$. Propomos aproximar a cauda da série por
$\sum_{k=3}^{6}g_k\,\mathrm{Li}_k(h)$. As constantes $g_k$ são determinadas por
um sistema linear que impõe a série exata até $h^5$ e o valor exato na elipse
degenerada. Como $\mathrm{Li}_3$ possui a mesma singularidade logarítmica
$(1-h)^2\ln(1-h)$ do perímetro exato, a aproximação é uniformemente boa em
todo o intervalo. O erro relativo máximo é $9{,}2\times10^{-9}$ e, para
$b/a\ge0{,}2$, $7{,}3\times10^{-11}$. Em todos os pontos testados ela é mais
precisa que a segunda aproximação de Ramanujan.
\end{abstract}

\section{A fórmula}

Sejam $a\ge b\ge0$ os semieixos e $h=\big(\frac{a-b}{a+b}\big)^2$. A aproximação proposta é
\begin{equation}\label{eq:formula}
\boxed{\;
P\approx\pi(a+b)\left[\,1+\frac h4+\frac{h^2}{64}
+\sum_{k=3}^{6}g_k\left(\mathrm{Li}_k(h)-h-\frac{h^2}{2^k}\right)\right]
\;}
\end{equation}
em que $\mathrm{Li}_k(h)=\sum_{n\ge1}h^n/n^k$ é o polilogaritmo de ordem $k$. As
constantes estão na Tabela~\ref{tab:g}.

\begin{table}[H]
\centering
\begin{tabular}{lr}
\toprule
constante & valor\\
\midrule
$g_3$ & $0.07953808367625804$\\
$g_4$ & $0.06062630887318473$\\
$g_5$ & $0.03505049186615928$\\
$g_6$ & $0.04933973528389241$\\
\bottomrule
\end{tabular}
\caption{Constantes da fórmula \eqref{eq:formula}. Elas são a solução exata do
sistema linear \eqref{eq:sistema}, e não fruto de ajuste numérico.}
\label{tab:g}
\end{table}

Os termos $-h-h^2/2^k$ retiram de cada polilogaritmo os seus dois primeiros
termos. Assim, os coeficientes de $h^0$, $h^1$ e $h^2$ da fórmula são exatamente
$1$, $\tfrac14$ e $\tfrac1{64}$.

\section{Dedução}

\subsection{A série exata e os números de Catalan}

O perímetro é $P=4aE(e)$, com $e^2=1-b^2/a^2$. A série de Gauss--Kummer dá
\begin{equation}\label{eq:serie}
P=\pi(a+b)\,f(h),\qquad f(h)=\sum_{n\ge0}a_nh^n,\qquad a_n=\binom{1/2}{n}^{\!2}.
\end{equation}
Nos extremos, $f(0)=1$ (círculo) e $f(1)=4/\pi$ (segmento, $P=4a$).

\begin{lema}\label{lema:catalan}
Seja $C_m=\frac{1}{m+1}\binom{2m}{m}$ o $m$-ésimo número de Catalan. Para $n\ge1$,
\[
\binom{1/2}{n}=\frac{(-1)^{n-1}\,C_{n-1}}{2^{2n-1}},
\qquad\text{logo}\qquad
a_n=\frac{C_{n-1}^{\,2}}{4^{\,2n-1}}.
\]
\end{lema}
\begin{proof}
Da definição,
$\binom{1/2}{n}=\frac{1}{n!}\prod_{j=0}^{n-1}\big(\tfrac12-j\big)
=\frac{(-1)^{n-1}}{2^n n!}\,1\cdot1\cdot3\cdot5\cdots(2n-3)$.
Usando $1\cdot3\cdots(2n-3)=\frac{(2n-2)!}{2^{n-1}(n-1)!}$, obtemos
$\binom{1/2}{n}=\frac{(-1)^{n-1}(2n-2)!}{2^{2n-1}\,n!\,(n-1)!}
=\frac{(-1)^{n-1}}{2^{2n-1}}\,C_{n-1}$.
\end{proof}

Assim,
$f(h)=1+\frac{h}{4}+\frac{h^2}{64}+\frac{h^3}{256}+\cdots$, e o
coeficiente de $h^n$ é o quadrado de $C_{n-1}/2^{2n-1}$, um objeto
combinatório: $C_{n-1}^2$ conta pares de caminhos de Dyck de comprimento
$2(n-1)$.

\subsection{Assintótica: o surgimento dos polilogaritmos}

A fórmula de Stirling dá a assintótica clássica dos números de Catalan,
$C_m\sim\dfrac{4^m}{\sqrt\pi\,m^{3/2}}$. Daí
\begin{equation}\label{eq:assint}
a_n=\frac{1}{4\pi n^3}\left(1+\frac{3}{4n}+O\!\left(\frac1{n^2}\right)\right),
\qquad n\to\infty.
\end{equation}
Numericamente, $4\pi n^3a_n=1.00075\ldots$ para $n=1000$, e
$(4\pi n^3a_n-1)\,n\to0.75$.

A expansão \eqref{eq:assint} diz que, para $n$ grande, $a_n$ é uma combinação
de potências $n^{-3},n^{-4},n^{-5},\dots$. A série cujos coeficientes são
exatamente $n^{-k}$ é o polilogaritmo:
\begin{equation}
\mathrm{Li}_k(h)=\sum_{n\ge1}\frac{h^n}{n^k},\qquad \mathrm{Li}_k(1)=\zeta(k),
\end{equation}
em que $\zeta$ é a função zeta de Riemann. Isso motiva a aproximação da
\emph{cauda} da série \eqref{eq:serie} por
\begin{equation}
\sum_{n\ge M+1}a_nh^n\;\approx\;\sum_{k=3}^{3+K}g_k\left(\mathrm{Li}_k(h)-\sum_{n=1}^{M}\frac{h^n}{n^k}\right),
\end{equation}
mantendo os termos $n\le M$ exatos. A fórmula \eqref{eq:formula} corresponde a
$M=2$ e $K=3$.

\subsection{A singularidade certa}

O perímetro exato não é analítico em $h=1$. Usando a expansão de $E(k)$ perto
de $k=1$, obtém-se
\begin{equation}
f(h)=\frac4\pi+\frac{1}{\pi}(h-1)-\frac{1}{8\pi}(1-h)^2\ln(1-h)+\cdots,
\end{equation}
e é por causa do termo $(1-h)^2\ln(1-h)$ que as aproximações racionais e
algébricas, como as de Ramanujan, perdem precisão perto do segmento.

\begin{proposicao}\label{prop:sing}
Perto de $h=1$,
$\mathrm{Li}_3(h)=\zeta(3)-\zeta(2)(1-h)-\tfrac12(1-h)^2\ln(1-h)+\cdots$, enquanto, para $k\ge4$,
$\mathrm{Li}_k$ só tem termos logarítmicos de ordem $(1-h)^{k-1}\ln(1-h)$.
\end{proposicao}
\begin{proof}
É a expansão clássica de $\mathrm{Li}_k(e^\mu)$ em torno de $\mu=0$:
$\mathrm{Li}_k(e^\mu)=\frac{\mu^{k-1}}{(k-1)!}\big[H_{k-1}-\ln(-\mu)\big]
+\sum_{j\ne k-1}\zeta(k-j)\frac{\mu^j}{j!}$, com $\mu=\ln h\approx-(1-h)$.
\end{proof}

Portanto a combinação \eqref{eq:formula} tem, perto do segmento, o termo
singular $-\frac{g_3}{2}(1-h)^2\ln(1-h)$. Pela assintótica \eqref{eq:assint},
deveríamos ter $g_3\approx\frac{1}{4\pi}$, o que daria exatamente o termo
$-\frac{1}{8\pi}(1-h)^2\ln(1-h)$ da função verdadeira. De fato,
\[
g_3=0.0795381\qquad\text{e}\qquad\frac1{4\pi}=0.0795775.
\]
A fórmula reproduz, quase sem erro, a singularidade correta. Esta é a razão de
sua precisão uniforme.

\subsection{O sistema linear}

As quatro constantes $g_3,\dots,g_6$ são definidas por quatro condições:
\begin{enumerate}
\item \textbf{Série exata em $h^3,h^4,h^5$.} O coeficiente de $h^n$ em
\eqref{eq:formula}, para $n\ge3$, é $\sum_k g_k\,n^{-k}$. Impomos que ele seja igual a $a_n$:
\item[] \begin{equation}\label{eq:sistema}
\sum_{k=3}^{6}\frac{g_k}{n^k}=a_n,\qquad n=3,4,5,
\end{equation}
com $a_3=\frac1{256}$, $a_4=\frac{25}{16384}$ e $a_5=\frac{49}{65536}$.
\item \textbf{Segmento exato.} Com $h=1$ e $\mathrm{Li}_k(1)=\zeta(k)$:
\begin{equation}
\sum_{k=3}^{6}g_k\left(\zeta(k)-1-2^{-k}\right)=\frac4\pi-1-\frac14-\frac1{64}.
\end{equation}
\end{enumerate}
A solução desse sistema $4\times4$ é a Tabela~\ref{tab:g}.

\subsection{Propriedades}

\begin{teorema}
A aproximação \eqref{eq:formula}:
\begin{enumerate}
\item é exata no círculo ($h=0$: $P=2\pi a$) e no segmento ($h=1$: $P=4a$);
\item coincide com a série exata \eqref{eq:serie} até o termo $h^5$, inclusive;
\item é contínua em $[0,1]$ e tem singularidade do tipo $(1-h)^2\ln(1-h)$ em $h=1$, com coeficiente
$-g_3/2$, a menos de $0{,}05\%$ do coeficiente exato $-1/(8\pi)$.
\end{enumerate}
\end{teorema}
\begin{proof}
(1) Em $h=0$, todos os termos com $h$ se anulam. Em $h=1$, vale por construção.
(2) Os coeficientes de $h^0,h^1,h^2$ são exatos por construção, e os de $h^3,h^4,h^5$ pelo
sistema \eqref{eq:sistema}.
(3) Segue da Proposição~\ref{prop:sing}, já que $\mathrm{Li}_k$ é contínua em $[0,1]$ para $k\ge2$.
\end{proof}

\section{Verificação numérica}

Comparamos com $P=4aE(e)$ calculado com 40 dígitos (\texttt{mpmath}), em 2000
valores de $h\in[0,1]$ e em pontos com $1-h=10^{-4},\dots,10^{-17}$.

\begin{table}[H]
\centering
\begin{tabular}{ccc}
\toprule
$b/a$ & erro relativo (nova) & erro relativo (Ramanujan II)\\
\midrule
0.9   & $2.8\times10^{-24}$ & $3.8\times10^{-18}$\\
0.7   & $5.7\times10^{-18}$ & $7.0\times10^{-13}$\\
0.5   & $1.3\times10^{-14}$ & $4.6\times10^{-10}$\\
0.3   & $4.8\times10^{-12}$ & $7.4\times10^{-8}$\\
0.2   & $7.4\times10^{-11}$ & $8.5\times10^{-7}$\\
0.1   & $1.1\times10^{-9}$  & $1.2\times10^{-5}$\\
0.05  & $4.5\times10^{-9}$  & $5.2\times10^{-5}$\\
0.01  & $9.0\times10^{-9}$  & $2.4\times10^{-4}$\\
0.001 & $2.4\times10^{-9}$  & $3.8\times10^{-4}$\\
0     & $0$ & $4.0\times10^{-4}$\\
\midrule
máximo ($b/a\ge0.2$) & $7.3\times10^{-11}$ & $8.5\times10^{-7}$\\
máximo geral & $9.2\times10^{-9}$ & $4.0\times10^{-4}$\\
\bottomrule
\end{tabular}
\caption{Erro relativo (em módulo) da fórmula \eqref{eq:formula} e de Ramanujan II.}
\end{table}

\begin{figure}[H]
\centering
\begin{tikzpicture}
\begin{axis}[
  width=0.85\textwidth, height=7cm,
  xlabel={$b/a$}, ylabel={erro relativo $|\,\cdot\,|$},
  ymode=log, ymin=1e-26, ymax=1e-2, xmin=0, xmax=1,
  legend pos=north east, grid=major]
\addplot[blue, thick] table[x=x, y=nova] {erros.dat};
\addlegendentry{nova (Eq.~\ref{eq:formula})}
\addplot[red, thick, dashed] table[x=x, y=ram2] {erros.dat};
\addlegendentry{Ramanujan II}
\end{axis}
\end{tikzpicture}
\caption{Erro relativo em função de $b/a$. Perto de $b/a=1$, os dois erros
ficam abaixo da precisão do cálculo, e as curvas são cortadas.}
\end{figure}

\section{Generalização e observações}

\begin{itemize}
\item \textbf{Família de fórmulas.} Mantendo $M$ termos exatos e usando $K+1$
polilogaritmos ($\mathrm{Li}_3,\dots,\mathrm{Li}_{3+K}$), o erro máximo cai
de forma sistemática. Por exemplo, $(M,K)=(0,2)$ dá $2{,}4\times10^{-5}$,
$(2,3)$ dá $9{,}2\times10^{-9}$ e $(4,4)$ dá $2{,}5\times10^{-11}$.
\item \textbf{Uso prático.} Os polilogaritmos estão disponíveis em
Mathematica (\texttt{PolyLog}), Maple, \texttt{mpmath} (\texttt{polylog}) e
na GSL. Não estão em calculadoras comuns.
\item \textbf{Ineditismo.} Em buscas na literatura (novembro de 2026), não
encontramos aproximações do perímetro da elipse baseadas em polilogaritmos ou
na assintótica dos números de Catalan. A busca não foi exaustiva e deve ser
complementada antes de qualquer publicação.
\end{itemize}

\end{document}
